% ============================================================
% APPENDIX B: EXTENDED EXPERIMENTS  (outside 15-page limit)
% Filled from Notion mirror + Part III PDF. Sample-size FLAG on dissociation n.
% ============================================================
\section{Extended experimental details}
\label{sec:extended}

\subsection{Full control-family seed table}
Table~\ref{tab:main} in the main text reports best-epoch (final-epoch in
parentheses). For completeness: condition means over available seeds are, best /
final, Phase~5 $3.606 / 3.567$; 6a $3.619 / 3.360$; 6b $3.601 / 3.415$; 6c
$3.544 / 3.294$. All values are EK100 action R@5 (\%), P01, signal-absent
evaluation, read from the per-run evaluation CSVs (five-decimal precision;
cross-checked against the tqdm logs). The 6b and 6c seed-44 runs were not
launched once the direction was closed, so those two conditions have two seeds.

\subsection{PEVA stage-1 and stage-2 statistics}
Stage-1 (average-error), $N=220$: mean relative gap $0.00055$; the real signal
beats the masked signal on $172/220$ clips ($78.2\%$); Wilcoxon
$p = 3.15\times10^{-19}$. Stage-2 (discrimination), $N=220$: top-1 error
identical for real and masked ($0.9545$; chance $0.0625$); McNemar both-correct
$210$, real-only $0$, masked-only $0$, both-wrong $10$; margin-$L_2$ Wilcoxon
$p = 2.46\times10^{-24}$. The two stages together give the detectable-but-not-
decision-relevant decoupling discussed in \S\ref{subsec:diagnostics}.

\subsection{Diagnostic raw outputs}
Output shift: $L_2(\text{real},\text{masked}) = 225.35$, $\lVert\text{real}\rVert
= 2272.18$, relative $L_2 = 0.0992$. Test~A (attention; uniform baseline
$0.001716$): gaze real $0.049080$ ($28.6\times$) vs.\ masked $0.049155$; hand
real $0.060709$ ($35.4\times$) vs.\ masked $0.061035$; real $>$ masked is false
for both. Test~B (direct probe; class-prior R@5 $= 0.1053$, chance $0.0041$),
per-target R@5 (vision, behavior, prior): action $0.1137 / 0.1216 / 0.1039$;
verb $0.4804 / 0.5941 / 0.6137$; noun $0.3176 / 0.2314 / 0.3000$. Note the verb
prior ($0.6137$) exceeds both arms: the behavioral advantage on verbs is relative
to vision, not to the prior.

\subsection{Verb/noun dissociation, all horizons}
Horizon sweep (fair-vision v3, common set), verb R@5 with verb prior $0.6749$:
3\,s behavior $0.5830\pm0.0108$, vision $0.5295\pm0.0305$, margin $+0.0535$
$[-0.0288, +0.0988]$; 5\,s $0.6145\pm0.0229$ vs.\ $0.5117\pm0.0118$, $+0.1029$
$[+0.0165, +0.1440]$; 10\,s $0.6159\pm0.0085$ vs.\ $0.4952\pm0.0216$, $+0.1207$
$[+0.0535, +0.1687]$. One-second row (separate subsample, prior $0.5922$):
behavior $0.5412\pm0.0096$, vision $0.5059\pm0.0169$, margin $+0.0353$
$[+0.0078, +0.1176]$. The earlier figure computed against the under-fit vision comparator is superseded by these values.
% FLAG: test-set n for the v3 common set -- Part III states 762; the reconstructed
% experiment log gives 243 for Option 2 v2 (with a known Notion 243->1. corruption
% now fixed). Confirm which n attaches to which row before stating a sample count.

\subsection{\texorpdfstring{$d \gg n$}{d >> n} comparator diagnostic}
In the signal-present HD-EPIC probe both arms interpolate their training data:
vision reaches training R@5 near unity while test R@5 sits at $0.50$--$0.62$; the
behavioral arm reaches training R@5 around $0.80$. A prior-mimicry guard voids any
run whose top-5 set collapses onto the class marginal (mean top-5 overlap
$\approx 1.0$). At roughly \num{1900} training examples the fair-vision pass-
condition fails at all four horizons, locating the constraint in sample size
rather than pathway design.